\graphicspath{{abstracts/byw/02-stadtmuseum-beispielstadt/img/}{img/}}

\begin{beitrag}{Sammlungsdokumentation am Stadtmuseum Beispielstadt: Eine WissKI-Migration mit kleinem Team}{%
Greta Gärtner\,\orcidlink{0000-0001-8901-2345} \\
\small \href{mailto:greta.gaertner@example.org}{greta.gaertner@example.org} \\
\small Friedrich-Alexander-Universität Erlangen-Nürnberg
}

\begin{abstract}
Das Projekt überführt die Sammlungsdokumentation des Stadtmuseums Beispielstadt (ca. 11.000 Objekte) aus einer veralteten Eigenentwicklung in eine WissKI-Umgebung mit einem kleinen Museumsteam; in den ersten zwölf Monaten wurde ein CIDOC-CRM-angelehntes Datenmodell entwickelt, das auch Ereignisse der Erwerbung, Restaurierung und Ausleihe abbildet, und die Pathbuilder-Konfiguration an einer 240-Objekte-Testinstanz finalisiert. Aktuell werden die ersten 1.500 Objekte des Hauptbestands migriert; zur Diskussion gestellt werden Organisationsmodelle für kleine Museumsteams, der produktive Umgang mit historisch gewachsenen Datenqualitätsproblemen sowie die Schnittstellenanbindung an einen regionalen Museumsverbund.
\end{abstract}

\section{Ausgangslage}
Das Stadtmuseum Beispielstadt verfügt über eine kulturhistorische Sammlung von etwa 11\,000 inventarisierten Objekten. Die bisherige Sammlungsdokumentation wurde über zwei Jahrzehnte in einer eigenentwickelten Datenbank gepflegt, deren technische Grundlage zunehmend veraltet ist und deren Inhalte für externe Recherchen nicht zugänglich sind. Im Rahmen einer dreijährigen Projektförderung wird die Sammlungsdokumentation nun in eine WissKI-Umgebung überführt; das Stadtmuseum kooperiert dafür mit dem WissKI-Konsortium an der FAU Erlangen-Nürnberg. Das Projekt steht aktuell an der Schwelle zwischen Modellierungs- und Migrationsphase und bringt sich im Format \emph{Bring your WissKI} mit konkreten Diskussionsfragen ein.

\section{Aktueller Stand}
In den ersten zwölf Monaten wurde die fachliche Modellierung der Sammlung in einem moderierten Prozess mit dem Museumsteam entwickelt. Diese Phase war erheblich aufwendiger als zunächst geplant, weil sich der größte Teil der bestehenden Datensätze als modellierungsrelevant herausstellte -- nicht nur die Objekteinträge, sondern auch die Vorgänge der Erwerbung, der Restaurierung und der Ausleihe sollten zukünftig im Datenmodell sichtbar bleiben. Die nun vorliegende Modellierung folgt einem an CIDOC CRM angelehnten Profil und ist im Detail dokumentiert.\footcite{gaertner2025-museumsmodell}

Parallel zur Modellierung wurde eine WissKI-Testinstanz aufgesetzt und mit einem ersten Sample von 240 Objekten befüllt. Die Erfahrungen aus dieser Erprobung sind in die finale Pathbuilder-Konfiguration eingeflossen, die seit etwa drei Monaten produktiv ist. Aktuell werden die ersten 1\,500 Objekte des Hauptbestands überführt; die vollständige Migration ist für die kommenden vierzehn Monate vorgesehen.

\section{Worüber wir diskutieren möchten}
Im Anwender*innentreffen möchten wir drei Themenfelder zur Diskussion stellen.

Erstens die \emph{Personal- und Kompetenzfrage}: Das Stadtmuseum verfügt über ein kleines Team -- eine wissenschaftliche Sammlungsleitung in Teilzeit, eine Sammlungsdokumentarin und eine halbe technische Stelle. Wie kann ein WissKI-Projekt in einer solchen Konstellation langfristig getragen werden, und welche Aufgaben können sinnvoll auf externe Partner ausgelagert werden, ohne dass die Kontrolle über die eigenen Daten verloren geht?

Zweitens die \emph{Datenqualitätsfrage}: Die bestehenden Datensätze sind über zwei Jahrzehnte gewachsen und enthalten erhebliche Inkonsistenzen. Wir würden gerne von anderen Museen lernen, wie sie die Datenqualität bei einer Migration verbessern, ohne dass die Migration in einer aufwendigen Bereinigung steckenbleibt; welche Datenqualitätsprobleme dürfen in die Zielumgebung übernommen werden, welche müssen vorab behoben werden?

Drittens die \emph{Schnittstellenfrage}: Das Stadtmuseum ist Mitglied eines regionalen Museumsverbunds, der eine gemeinsame Recherche-Plattform betreibt. Wir wünschen uns einen Austausch mit Häusern, die ihre WissKI-Bestände in vergleichbaren Verbund-Plattformen sichtbar machen, und sind interessiert an konkreten Erfahrungen mit den dafür benötigten Datenstandards und Schnittstellen.\footcite{gaertner2025-verbund}

\section{Vorgehensweise im Format \emph{Bring your WissKI}}
Wir werden in einem kurzen, etwa zwanzigminütigen Aufschlag die zentralen Designentscheidungen unseres Projekts vorstellen und drei konkrete Datenbeispiele zeigen, an denen sich die genannten Diskussionsfragen festmachen lassen. Anschließend möchten wir in eine offene Gesprächsrunde übergehen, in der wir gemeinsam mit den Teilnehmer*innen erörtern, wie sich die Themen in vergleichbaren Häusern darstellen.

\section{Was wir mitbringen, was wir suchen}
Wir bringen mit: eine produktiv laufende, dokumentierte WissKI-Konfiguration, ein an CIDOC CRM angelehntes Datenmodell mit einer ausführlichen Begründung der Designentscheidungen, und konkrete Erfahrungen mit der Migration aus einer historisch gewachsenen Datenbank. Wir suchen: einen Austausch mit Museen ähnlicher Größenordnung, die vergleichbare Migrationsvorhaben durchführen oder bereits abgeschlossen haben, und sind besonders interessiert an Erfahrungen mit kleinem Team und an Wegen zur langfristigen institutionellen Verstetigung.

\end{beitrag}

\section*{Kurzbiografie(n)}

\subsection*{Greta Gärtner}
Greta Gärtner ist Postdoktorandin an der FAU Erlangen-Nürnberg und Mitglied des WissKI-Konsortiums. Ihre Forschungsschwerpunkte liegen in der Anwendung semantischer Datenmodelle auf museale Sammlungen mit kleinem und mittelgroßem Bestand. Sie hat in mehreren Migrationsprojekten als wissenschaftliche Begleitung mitgewirkt und beschäftigt sich aktuell mit der Schnittstellenanbindung von WissKI-Beständen an Museums-Verbund-Plattformen.
